
@@M0@@
\item \uline{Le
conoyau} @@M1@@ \uline{de} @@M2@@ \uline{est un} @@M3@@\uline{-module localement libre.}
\item \uline{L'homomorphisme} @@M4@@ \uline{est surjectif.}
\end{enumerate}
\uline{Alors le morphisme canonique} @@M5@@ \uline{est représentable par un morphisme affine de présentation finie~; plus précisément, on a une suite exacte canonique de faisceaux} f.p.p.f.,
@@M6@@
\uline{et, si} @@M7@@\uline{, le monomorphisme} @@M8@@ \uline{est représentable par une immersion fermée et} @@M9@@ \uline{est un torseur sur} @@M10@@ \uline{sous le groupe} @@M11@@ \uline{(c'est-à-dire, la suite de faisceaux} f.p.p.f.
@@M12@@
\uline{est exacte).}

Considérons pour un @@M13@@-schéma @@M14@@ les deux diagrammes suivants à lignes exactes :
@@M15@@
@@M16@@

On fera attention que le carré de gauche de ce diagramme \uline{n'est pas nécessairement commutatif}. Cependant on a~:
@@M17@@
